Aos meus pais, Veronico e Ana Paula, pelo amor incondicional, pelo esforço diário e por serem a base que me sustentou em cada etapa desta caminhada. Sem o apoio e a confiança de vocês, nada disso seria possível.

À minha namorada, Lunna, pelo companheirismo, pelo carinho e pela paciência de sempre. Obrigado por caminhar ao meu lado, celebrar cada pequena conquista e me dar forças quando o cansaço batia.

Ao meu padrinho José \textit{(in memoriam)}, cuja lembrança permanece viva comigo, e à minha madrinha Dila, pelo carinho e cuidado constantes. Agradeço também aos meus tios e primos, que acompanharam minha trajetória de perto e sempre estiveram prontos para ajudar.

Ao meu orientador, Prof. Dr. Matheus Fernandes de Araújo Silva, pela parceria, pela confiança e pela orientação ao longo deste trabalho — e, principalmente, por aguentar com bom humor a inversão de papéis em que o orientando é quem vivia cobrando o orientador.

Aos meus companheiros da Competição de Pontes com Palitos de Picolé (CP3), projeto do qual tenho orgulho de fazer parte como colaborador, pela parceria nos bastidores e pela vivência que deu origem a este trabalho. Um agradecimento especial à nossa ``média culpula'', pela resenha, pela amizade e por estarem sempre juntos fazendo tudo acontecer.

Aos meus amigos e colegas de jornada, que dividiram comigo o peso dos momentos difíceis e tornaram a rotina acadêmica muito mais leve. Esta graduação teve seus altos e baixos, mas contar com a amizade de vocês fez toda a diferença para que eu chegasse até o fim.

Por fim, à Universidade Federal Rural do Semi-Árido (UFERSA), campus Pau dos Ferros, e a todos os professores do curso de Tecnologia da Informação que contribuíram para a minha formação acadêmica e pessoal.